\documentclass[10pt, aspectratio=169]{beamer}
\usepackage{../beamer_theme_408}

\title{408考研零基础 --- 数据结构}
\subtitle{1-8 数组与特殊矩阵}
\author{}
\date{}


\begin{document}


\begin{frame}{本节目标}
    \begin{enumerate}
        \item 理解数组的两种存储方式及地址计算方法
        \item 掌握对称矩阵、三角矩阵、三对角矩阵的压缩存储与地址映射
        \item 理解稀疏矩阵的三元组存储方法
    \end{enumerate}
\end{frame}

\begin{frame}{数组的存储结构}
    \textbf{二维数组 $A[m][n]$}
    \vspace{0.3em}
    \textbf{行优先存储}
    \[
    \text{LOC}(a_{ij}) = \text{LOC}(a_{00}) + (i \times n + j) \times L
    \]
    \vspace{0.3em}
    \textbf{列优先存储}
    \[
    \text{LOC}(a_{ij}) = \text{LOC}(a_{00}) + (j \times m + i) \times L
    \]
    \vspace{0.3em}
    \begin{block}{C 语言}
        采用行优先存储
    \end{block}
\end{frame}

\begin{frame}{对称矩阵的压缩}
    \textbf{$n$ 阶对称矩阵}：$a_{ij} = a_{ji}$
    \vspace{0.3em}
    \textbf{存储下三角 $+$ 主对角线}
    \[
    \text{元素个数} = \frac{n(n+1)}{2}
    \]
    \vspace{0.3em}
    \textbf{地址映射}（$i \ge j$）
    \[
    \text{下标} = \frac{i(i+1)}{2} + j
    \]
    \vspace{0.3em}
    当 $i < j$ 时，访问对称元素 $a_{ji}$
\end{frame}

\begin{frame}{三角矩阵}
    \textbf{下三角矩阵}
    \begin{itemize}
        \item 上三角区域全部为常数 $c$
        \item 存储 $n(n+1)/2 + 1$ 个元素
        \item 地址映射与对称矩阵相同（$i \ge j$）
        \item $i < j$ 时返回常数 $c$
    \end{itemize}
    \vspace{0.3em}
    \textbf{上三角矩阵}
    \begin{itemize}
        \item 按列优先存储上三角区域
    \end{itemize}
\end{frame}

\begin{frame}{三对角矩阵}
    \textbf{非零元素在三条对角线上}
    \[
    |i - j| \le 1
    \]
    \vspace{0.3em}
    \textbf{按行优先存储}
    \[
    \text{下标} = 2i + j
    \]
    \vspace{0.3em}
    \textbf{举例（$5 \times 5$）}
    \[
    a_{00}, a_{01}, a_{10}, a_{11}, a_{12}, a_{21}, a_{22}, a_{23}, a_{32}, a_{33}, a_{34}, a_{43}, a_{44}
    \]
\end{frame}

\begin{frame}{稀疏矩阵 --- 三元组表示}
    \textbf{只存储非零元素}
    \[
    (row, col, value)
    \]
    \vspace{0.3em}
    \begin{center}
    \begin{tabular}{|c|c|c|}
        \hline
        row & col & value \\
        \hline
        0 & 0 & 3 \\
        \hline
        1 & 2 & 5 \\
        \hline
        2 & 0 & 7 \\
        \hline
        \vdots & \vdots & \vdots \\
        \hline
    \end{tabular}
    \end{center}
    \vspace{0.3em}
    \textbf{缺点}：失去随机访问能力
\end{frame}

\begin{frame}{稀疏矩阵 --- 十字链表}
    \textbf{每个非零元素结点}
    \begin{itemize}
        \item 行号、列号、值
        \item 指向同一行下一个非零元素的指针（右指针）
        \item 指向同一列下一个非零元素的指针（下指针）
    \end{itemize}
    \vspace{0.5em}
    优点：支持快速存取任意位置的元素
\end{frame}

\begin{frame}{本节总结}
    \begin{enumerate}
        \item 数组存储：行优先和列优先两种方式
        \item 对称矩阵：存储 $n(n+1)/2$ 个元素
        \item 三角矩阵：下三角/上三角 $+$ 常数
        \item 三对角矩阵：存储三条对角线
        \item 稀疏矩阵：三元组或十字链表
    \end{enumerate}
\end{frame}

\end{document}
